\section{The exact finite complement}
\label{sec:finite}

Only diameters $2\le D\le15$ remain. If there is no interior
coordinate, Lemma~\ref{lem:affine} already gives the two-symbol bound.
Otherwise some $t\in E'\cap\{1,\ldots,D-1\}$ exists, and
\eqref{eq:secant-bound} implies
\[
 |t-r|\le\frac{2D}{t(D-t)}
 \le\frac{2D}{D-1}\le4.
\]
Consequently $-3\le r\le D+3$. Combining this with the endpoint
conditions gives the complete finite ranges
\begin{equation}
 \begin{gathered}
 2\le D\le15,\qquad -3\le r\le D+3,\qquad 0\le A\le2D,\\
 q\in\{-2,-1,0,1,2\},\qquad 0\le A+qD\le2D.
 \end{gathered}
 \label{eq:finite-ranges}
\end{equation}
This bounds the normalized secant parameters after the global
reduction; it does not impose a cutoff on $a,b,c$.

For each tuple define
\begin{equation}
 S=\{t\in\Z:0\le t\le D,\ |t(t-D)(t-r)|\le2D\}
 \label{eq:small-symbols}
\end{equation}
and the partial map on $S^2$
\begin{equation}
 F(x,y)=(y,A+qy+y(y-D)(y-r)-x).
 \label{eq:small-graph}
\end{equation}
An image leaving $S^2$ has no outgoing graph edge. Every actual
periodic point of a global-extrema tuple lies on a cycle of this
graph. Conversely, every graph cycle is an actual polynomial orbit.
Retain a tuple only if the union of its graph cycles uses both $0$
and $D$. This is a necessary condition for an actual global-extrema
tuple, and hence cannot discard the periodic set of any map under
consideration.

\begin{proposition}[Finite arithmetic certificate]
\label{prop:finite}
The complete computation for \eqref{eq:finite-ranges} gives
Table~\ref{tab:small}. Every retained graph cycle uses at most
three symbols. Its least period belongs to $\{1,2,3,4,6\}$.
The largest total number of graph-periodic vertices is eleven,
attained at the single tuple
\begin{equation}
 (D,r,A,q)=(4,2,6,-1).
 \label{eq:max-tuple}
\end{equation}
The cycles of this tuple are
\begin{equation}
 (2),\quad (0,3),\quad (1,4),\quad (0,2,4),\quad (0,4,2).
 \label{eq:max-words}
\end{equation}
\end{proposition}

\begin{table}[htbp]
\centering\small
\begin{tabular}{@{}r r r r c@{}}
\toprule
$D$ & Compatible tuples & Both graph ends & Maximum points & Periods\\
\midrule
$2$  & $117$  & $51$  & $9$  & $1,2,3,4,6$\\
$3$  & $170$  & $56$  & $9$  & $1,2,3,4$\\
$4$  & $231$  & $63$  & $11$ & $1,2,3,4,6$\\
$5$  & $300$  & $68$  & $8$  & $1,2,3,4$\\
$6$  & $377$  & $75$  & $9$  & $1,2,3,4,6$\\
$7$  & $462$  & $78$  & $8$  & $1,2,3,4$\\
$8$  & $555$  & $85$  & $9$  & $1,2,3,4,6$\\
$9$  & $656$  & $90$  & $8$  & $1,2,3,4$\\
$10$ & $765$  & $97$  & $9$  & $1,2,3,4,6$\\
$11$ & $882$  & $102$ & $8$  & $1,2,3,4$\\
$12$ & $1007$ & $109$ & $9$  & $1,2,3,4,6$\\
$13$ & $1140$ & $114$ & $8$  & $1,2,3,4$\\
$14$ & $1281$ & $121$ & $9$  & $1,2,3,4,6$\\
$15$ & $1430$ & $126$ & $8$  & $1,2,3,4$\\
\bottomrule
\end{tabular}
\caption{The full small-diameter complement: 9,373 tuples satisfy the
endpoint value inequalities and 1,235 use both endpoints among graph
cycles. Point maxima and period lists are taken after this necessary
filter. In every row the maximum number of symbols per cycle is three.}
\label{tab:small}
\end{table}

\begin{proof}
For each tuple in \eqref{eq:finite-ranges}, form \eqref{eq:small-symbols},
evaluate \eqref{eq:small-graph} exactly, and extract every directed
cycle. Retain precisely those tuples passing the endpoint filter.
Sum the lengths of the distinct cycles to count periodic vertices.
Appendix~\ref{app:certificate} supplies the complete integer
pseudocode, including the extraction algorithm, the filter, the
maximizer update and the period and symbol statistics. Its exact
evaluation is Table~\ref{tab:small}, \eqref{eq:max-tuple} and
\eqref{eq:max-words}.

The first count in each row can also be checked directly. There are
$D+7$ choices for $r$. At $q=-2,-1,0,1,2$, the numbers of allowed
$A$ values are respectively $1,D+1,2D+1,D+1,1$. Their sum is
$4D+5$, giving $(D+7)(4D+5)$ tuples. The graph algorithm, rather
than this input count, establishes every cycle statistic.

All arithmetic decisions are integer equalities or inequalities.
The graph is finite, and the extraction lemma in the appendix proves
that it records every cycle exactly once and records its least period.
No assumed upper bound eleven enters the computation of the maximum;
the maximum and all tuples attaining it are accumulated from the graphs.
The reduction preceding \eqref{eq:finite-ranges} proves that this
finite computation covers the entire remaining family.
\end{proof}

Together, Lemmas~\ref{lem:integrality}--\ref{lem:affine} and
Propositions~\ref{prop:large-computation}--\ref{prop:finite} prove
the eleven-point bound, the possible period list and the per-cycle
three-symbol restriction. Their roles are different: the analytic
arguments prove the finite reduction, while the exact certificate
evaluates its explicitly bounded graphs. The next section uses
these established restrictions to derive the full coefficient
conditions without any bound on the original coefficients.
